\chapter{Aprendizaje estadístico y sobreajuste}
\label{cap:aprendizaje}

\section{Formulación del problema}

Sea $(\bm x, y)$ un par aleatorio con distribución desconocida $\mathcal{D}$, donde $\bm x\in\R^d$ son las características espectrales de una ventana e $y\in\{0,1\}$ indica si hay crisis. Se busca una función $h_\theta:\R^d\to[0,1]$ que minimice el \emph{riesgo}
\begin{equation}
R(\theta) = \E_{(\bm x,y)\sim\mathcal D}\big[\ell(h_\theta(\bm x), y)\big],
\end{equation}
pero como $\mathcal D$ es desconocida se minimiza el \emph{riesgo empírico} $\hat R(\theta) = \frac1n\sum_i \ell(h_\theta(\bm x_i), y_i)$ sobre una muestra de entrenamiento.

\section{Entropía cruzada y máxima verosimilitud}

Si se modela $P(y=1\mid\bm x) = \sigma(z_\theta(\bm x))$, con $\sigma(z) = 1/(1+e^{-z})$, la log-verosimilitud negativa de la muestra es
\begin{equation}
-\log\mathcal{L}(\theta) = \sum_i \Big[-y_i\log\sigma(z_i) - (1-y_i)\log(1-\sigma(z_i))\Big] = \sum_i\Big[\log(1+e^{z_i}) - y_i z_i\Big],
\end{equation}
que es la \emph{entropía cruzada binaria} \parencite[cap.~5]{bishop2024}. Minimizarla equivale a estimar $\theta$ por máxima verosimilitud.

\subsection{Desbalance de clases}

Con una prevalencia de crisis $\pi\approx 0{,}01$, un clasificador que siempre responde ``normal'' tiene una exactitud del 99~\% y sensibilidad nula. Para compensar se pondera la pérdida: $\ell_i = c_{y_i}\,[\log(1+e^{z_i}) - y_i z_i]$ con $c_1/c_0 = (1-\pi)/\pi$. Esto equivale a entrenar como si las clases estuvieran balanceadas; el umbral de decisión se ajusta después (capítulo~\ref{cap:evaluacion}).

\pendiente{teoría de decisión bayesiana y riesgo con costos asimétricos, \cite[cap.~2]{duda2001}.}

\section{Sesgo, varianza y sobreajuste}

Para la pérdida cuadrática, el error esperado de un estimador $\hat h$ entrenado sobre una muestra aleatoria se descompone en
\begin{equation}
\E\big[(y-\hat h(\bm x))^2\big] = \underbrace{\big(\E[\hat h(\bm x)] - h^\star(\bm x)\big)^2}_{\text{sesgo}^2} + \underbrace{\operatorname{Var}\big[\hat h(\bm x)\big]}_{\text{varianza}} + \underbrace{\sigma^2}_{\text{ruido}}.
\end{equation}
Un modelo con mucha capacidad reduce el sesgo pero aumenta la varianza: se ajusta a particularidades de la muestra de entrenamiento (\emph{sobreajuste}) y generaliza mal.

\pendiente{demostración de la descomposición; figura conceptual.}

\section{Regularización}

\begin{itemize}
  \item \textbf{Penalización $L_2$:} minimizar $\hat R(\theta) + \frac\mu2\|\theta\|^2$. Por dualidad de Lagrange es equivalente a minimizar $\hat R$ sujeto a $\|\theta\|^2\le r^2$ para algún $r$ que depende de $\mu$ \parencite[sec.~7.2]{goodfellow2016}. En AdamW se aplica como decaimiento de pesos desacoplado \parencite{loshchilov2019}.
  \item \textbf{Dropout:} en cada paso se anula aleatoriamente cada unidad oculta con probabilidad $q$, lo que equivale aproximadamente a promediar un ensamble de subredes.
  \item \textbf{Early stopping:} se detiene el entrenamiento cuando el desempeño en validación deja de mejorar. Para modelos lineales con pérdida cuadrática es equivalente a una regularización $L_2$ \parencite[sec.~7.8]{goodfellow2016}.
\end{itemize}

\section{Fuga de información y partición por paciente}

Las ventanas consecutivas de un mismo registro son casi idénticas. Si se reparten al azar entre entrenamiento y prueba, el modelo puede ``reconocer al paciente'' en lugar de aprender a detectar crisis, y la métrica de prueba queda inflada. Por eso todas las ventanas de un mismo sujeto van al mismo conjunto. En CHB-MIT, además, los casos \texttt{chb01} y \texttt{chb21} son la misma paciente y se tratan como un único sujeto.

\pendiente{comparación partición por sujeto vs. aleatoria (sección 11 del notebook), figura \texttt{sobreajuste\_curvas}.}
